\documentclass[11pt,a4paper]{article}
\usepackage[margin=24mm]{geometry}
\usepackage[T1]{fontenc}
\usepackage{lmodern}
\usepackage{amsmath,amssymb,amsthm,mathtools}
\usepackage[hidelinks]{hyperref}
\usepackage{bookmark}
\setlength{\emergencystretch}{2em}
\newtheorem{proposition}{Proposition}
\newcommand{\Q}{\mathbb Q}
\newcommand{\R}{\mathbb R}
\newcommand{\OO}{\mathcal O}
\DeclareMathOperator{\coeff}{coeff}
\title{Valuative Decay of Normalized Pushforwards}
\author{}
\date{}
\hypersetup{pdftitle={Valuative Decay of Normalized Pushforwards},
  pdfsubject={A proof for effective Weil real divisors}}
\begin{document}
\maketitle
\vspace{-2.2em}

Let $X$ be a normal projective variety over an arbitrary field $k$.
For a Weil $\R$-divisor $B$ and a prime divisor
$F$ on a normal projective birational model $p:Y\to X$, write
$v_F(B)=\coeff_F(p^*B)$, where $p^*$ denotes the asymptotic pullback.
We use its monotonicity and its agreement with the usual pullback on
$\Q$-Cartier divisors. All divisor inequalities are coefficientwise;
intersections of Cartier divisors with Weil divisors are understood
as intersection numbers of cycles.

\begin{proposition}
Let $f:X\to X$ be a finite surjective int-amplified endomorphism.
For every effective Weil $\R$-divisor $D$ and every divisorial valuation
$v_F$ over $X$, one has
\[
 \lim_{i\to\infty}
 v_F\!\left(\frac{(f^i)_*D}{\deg(f^i)}\right)=0.
\]
\end{proposition}

\begin{proof}
Choose an ample Cartier divisor $H$ and a rational number $\lambda>1$
such that $f^*H-\lambda H$ is ample.

\smallskip
\noindent\textit{Step 1. Construct effective $\Q$-Cartier upper bounds.}
Choose an effective Cartier divisor $C\ge D$ and an integer $m>0$
such that $mH-C$ is ample. Such a $C$ exists: for $a\gg0$, a nonzero
section of $\OO_X(aH-\lceil D\rceil)$ gives a divisor $C\in|aH|$
dominating $D$. For $i\ge1$, the identity
\[
 (f^i)^*H-\lambda^iH
 =\sum_{j=0}^{i-1}\lambda^{i-1-j}(f^j)^*(f^*H-\lambda H)
\]
shows that $(f^i)^*H-\lambda^iH$ is ample, since finite pullbacks
preserve ampleness. Hence
\[
 m\lambda^{-i}(f^i)^*H-C
 =(mH-C)+m\lambda^{-i}\bigl((f^i)^*H-\lambda^iH\bigr)
\]
is ample. Take an effective $\Q$-Cartier divisor
$G_i\sim_{\Q}m\lambda^{-i}(f^i)^*H-C$, and set
\[
 M_i:=\frac{(f^i)_*(C+G_i)}{\deg(f^i)}.
\]
Since $C+G_i\ge D$, we obtain
\begin{equation}\label{eq:upper}
 0\le \frac{(f^i)_*D}{\deg(f^i)}\le M_i\sim_{\Q}m\lambda^{-i}H.
\end{equation}
Here pushforward preserves $\Q$-linear equivalence by the norm formula,
and $(f^i)_*(f^i)^*H=\deg(f^i)H$. In particular, $M_i$ is
$\Q$-Cartier.

\smallskip
\noindent\textit{Step 2. Estimate the coefficient on a fixed model.}
Fix a normal projective birational model $p:Y\to X$ carrying $F$,
and choose an ample Cartier divisor $A$ on $Y$. Put $n=\dim X$ and
$a_{F,i}=v_F((f^i)_*D/\deg(f^i))$. By monotonicity and
\eqref{eq:upper},
\[
 0\le a_{F,i}F\le p^*M_i.
\]
Intersecting with $A^{n-1}$ gives
\[
 \begin{aligned}
 0\le a_{F,i}(F\cdot A^{n-1})
 &\le p^*M_i\cdot A^{n-1}\\
 &=m\lambda^{-i}(p^*H\cdot A^{n-1}).
 \end{aligned}
\]
Since $A$ is ample, the prime cycle $F$ has positive degree
$F\cdot A^{n-1}>0$. It follows that
\[
 0\le v_F\!\left(\frac{(f^i)_*D}{\deg(f^i)}\right)=a_{F,i}
 \le m\lambda^{-i}
 \frac{p^*H\cdot A^{n-1}}{F\cdot A^{n-1}}
 \longrightarrow0.
\]
\end{proof}

\end{document}
